\documentclass[11pt,letterpaper]{article}
\usepackage[utf8]{inputenc}
\usepackage[T1]{fontenc}
\usepackage[spanish,es-tabla]{babel}
\usepackage[margin=2.5cm]{geometry}
\usepackage{xcolor}
\usepackage{hyperref}
\usepackage{enumitem}

\hypersetup{
    colorlinks=true,
    linkcolor=blue!50!black,
    urlcolor=blue!50!black,
    pdftitle={Anexo - Fichas bibliograficas con extracto citado},
    pdfauthor={Equipo ACC}
}

\title{\textbf{Anexo --- Fichas bibliográficas}\\
\large Extracto exacto usado de cada fuente en el informe de la Práctica Experimental Unidad 5}
\author{Equipo ACC}
\date{Agosto 2026}

\newcommand{\ficha}[7]{%
\subsection*{#1}
\noindent\textbf{Autores:} #2 \\
\textbf{Año / Venue:} #3 \\
\textbf{DOI:} \url{#4} \\
\textbf{Enlace directo (acceso abierto, arXiv):} \url{#5}\\[0.5em]
\noindent\textbf{Fragmento usado en el informe} (parágrafo #6 --- resaltado abajo):\\[0.3em]
\noindent\colorbox{yellow!60}{\parbox{0.94\linewidth}{\vspace{0.3em}#7\vspace{0.3em}}}\\[0.3em]
\noindent\hrulefill
}

\begin{document}
\maketitle
\thispagestyle{empty}

\noindent Este anexo documenta, para cada fuente académica citada en el informe principal, exactamente qué idea o hallazgo se usó y en qué parte del texto. Todas las fuentes son artículos de acceso abierto en arXiv.org (preprints públicos, sin restricción de copyright para su lectura y cita); no se utilizó ningún repositorio de acceso no autorizado a contenido con derechos de autor para obtenerlas. El resaltado amarillo marca el fragmento --- en este caso una paráfrasis fiel, no una copia textual extensa del PDF original --- que respalda la afirmación correspondiente del informe.

\vspace{1em}

\ficha{1. Analyzing the Effects of CI/CD on Open Source Repositories in GitHub and GitLab}
{Jeffrey Fairbanks, Akshharaa Tharigonda, Nasir U. Eisty (Boise State University)}
{2023, arXiv (cs.SE)}
{https://doi.org/10.48550/arXiv.2303.16393}
{https://arxiv.org/abs/2303.16393}
{Sección 2 --- Metodología de verificación}
{El estudio analizó más de 12{,}000 repositorios de GitHub y GitLab y midió el efecto de la adopción de CI/CD sobre la velocidad de commits y el número de issues reportados, mostrando que el cumplimiento de CI/CD debe verificarse sobre artefactos reales del repositorio (workflows, historial de builds) y no solo declararse.}

\ficha{2. Efficacy of Static Analysis Tools for Software Defect Detection on Open-Source Projects}
{Jones Yeboah, Saheed Popoola}
{2024, arXiv (cs.SE)}
{https://doi.org/10.48550/arXiv.2405.12333}
{https://arxiv.org/abs/2405.12333}
{Sección 3.3 --- Análisis estático de código (criterio c)}
{Comparando SonarQube, PMD, Checkstyle y FindBugs sobre proyectos en Java, C++ y Python, los autores encontraron que SonarQube obtuvo mejor precisión y recall que las otras tres herramientas en la detección de defectos, lo que respalda recomendar SonarCloud como opción unificada para cerrar la brecha de análisis estático del pipeline.}

\ficha{3. Metric Criticality Identification for Cloud Microservices}
{Akanksha Singal, Divya Pathak, Kaustabha Ray, Felix George, Mudit Verma, Pratibha Moogi (IBM Research India)}
{2025, arXiv (cs.DC)}
{https://doi.org/10.48550/arXiv.2501.03547}
{https://arxiv.org/abs/2501.03547}
{Sección 4.2 --- Métricas de negocio en \texttt{/metrics} (criterio b)}
{Los autores proponen KIMetrix, un método que usa entropía e información mutua para identificar el subconjunto mínimo de métricas verdaderamente relevante para operar un sistema de microservicios, en lugar de exponer métricas genéricas sin curar --- el mismo criterio de métricas específicas del dominio (reintentos de transacción, escalamiento de tickets) aplicado en el proyecto evaluado.}

\ficha{4. LO2: Microservice API Anomaly Dataset of Logs and Metrics}
{Alexander Bakhtin, Jesse Nyyssölä, Yuqing Wang, Noman Ahmad, Ke Ping, Matteo Esposito, Mika Mäntylä, Davide Taibi (Universidad de Oulu / Universidad de Helsinki)}
{2025, arXiv (cs.SE)}
{https://doi.org/10.48550/arXiv.2504.12067}
{https://arxiv.org/abs/2504.12067}
{Sección 4.1 --- Logging estructurado (criterio a)}
{Los autores publican un dataset real de producción con cerca de 657{,}000 archivos de log y más de 45 millones de archivos de métricas, y muestran que la detección de anomalías en microservicios depende de poder fusionar logs y métricas de forma consistente --- lo que exige una estructura de log parseable en todos los servicios, no solo en algunos.}

\ficha{5. Informed and Assessable Observability Design Decisions in Cloud-native Microservice Applications}
{Maria C. Borges, Joshua Bauer, Sebastian Werner, Michael Gebauer, Stefan Tai (TU Berlin) --- ICSA 2024}
{2024, IEEE ICSA / arXiv (cs.SE)}
{https://doi.org/10.48550/arXiv.2403.00633}
{https://arxiv.org/abs/2403.00633}
{Sección 4.3 --- Dashboard de observabilidad (criterio c)}
{El trabajo argumenta que las decisiones de diseño de observabilidad (qué se instrumenta, dónde se visualiza) deben ser explícitas y evaluables mediante experimentos, en lugar de asumirse por la sola presencia de herramientas; esto respalda documentar explícitamente la desviación de plataforma (Grafana self-hosted vs. Grafana Cloud) en vez de dejarla implícita.}

\end{document}
